\documentclass[11pt]{article}
\usepackage[margin=1in]{geometry}
\usepackage{enumitem}
\usepackage{titlesec}
\usepackage{hyperref}
% AUTO-GENERATED by reports/generate_macros.py from reports/data/phase_quantities.json
% Do not hand-edit -- update the JSON and regenerate.

% --- phase1 ---
\newcommand{\PhaseOneHaSpearman}{0.115}
\newcommand{\PhaseOneHaCiLo}{0.102}
\newcommand{\PhaseOneHaCiHi}{0.127}
\newcommand{\PhaseOneHaN}{23983}
\newcommand{\PhaseOneKmSpearman}{0.457}
\newcommand{\PhaseOneKmCiLo}{0.434}
\newcommand{\PhaseOneKmCiHi}{0.48}
\newcommand{\PhaseOneKmN}{4722}
\newcommand{\PhaseOnePsSpearman}{0.295}
\newcommand{\PhaseOnePsCiLo}{0.269}
\newcommand{\PhaseOnePsCiHi}{0.321}
\newcommand{\PhaseOnePsN}{4611}
\newcommand{\PhaseOneXhSpearman}{0.342}
\newcommand{\PhaseOneXhCiLo}{0.307}
\newcommand{\PhaseOneXhCiHi}{0.378}
\newcommand{\PhaseOneXhN}{2502}
\newcommand{\PhaseOneZuSpearman}{0.247}
\newcommand{\PhaseOneZuCiLo}{0.212}
\newcommand{\PhaseOneZuCiHi}{0.283}
\newcommand{\PhaseOneZuN}{2968}
\newcommand{\PhaseOneEndeBenchmarkSpearman}{0.57}
\newcommand{\PhaseOneFlagThreshold}{0.3}
\newcommand{\PhaseOneNPairsFlagged}{3}
\newcommand{\PhaseOneNPairsTotal}{5}

% --- phase2 ---
\newcommand{\PhaseTwoOovPooledPearson}{-0.33}
\newcommand{\PhaseTwoOovPooledAuc}{0.513}
\newcommand{\PhaseTwoOovPerlangLo}{-0.045}
\newcommand{\PhaseTwoOovPerlangHi}{0.021}
\newcommand{\PhaseTwoSentlenPooledPearson}{0.19}
\newcommand{\PhaseTwoSentlenPooledAuc}{0.503}
\newcommand{\PhaseTwoIdiomPooledPearson}{-0.07}
\newcommand{\PhaseTwoIdiomPooledAuc}{0.517}
\newcommand{\PhaseTwoModelUniversalAuc}{0.513}
\newcommand{\PhaseTwoModelUniversalN}{38786}
\newcommand{\PhaseTwoModelIdiomAuc}{0.52}
\newcommand{\PhaseTwoModelFullAuc}{0.564}
\newcommand{\PhaseTwoModelFullN}{10812}
\newcommand{\PhaseTwoAucChance}{0.5}

% --- phase3 ---
\newcommand{\PhaseThreeModel}{Aya-101}
\newcommand{\PhaseThreeNSentencesPerPair}{100}
\newcommand{\PhaseThreeEndeBaselineCometkiwi}{0.804}
\newcommand{\PhaseThreeEndeBaselineComet}{0.821}
\newcommand{\PhaseThreeEndeBaselineChrf}{53.3}
\newcommand{\PhaseThreeEndeCotCometkiwi}{0.766}
\newcommand{\PhaseThreeEndeCotComet}{0.782}
\newcommand{\PhaseThreeEndeCotChrf}{50.4}
\newcommand{\PhaseThreeEndeToolCometkiwi}{0.819}
\newcommand{\PhaseThreeEndeToolComet}{0.843}
\newcommand{\PhaseThreeEndeToolChrf}{59.5}
\newcommand{\PhaseThreeEndeCottoolCometkiwi}{0.72}
\newcommand{\PhaseThreeEndeCottoolComet}{0.733}
\newcommand{\PhaseThreeEndeCottoolChrf}{48.2}
\newcommand{\PhaseThreeEnhaBaselineCometkiwi}{0.609}
\newcommand{\PhaseThreeEnhaBaselineComet}{0.762}
\newcommand{\PhaseThreeEnhaBaselineChrf}{42.2}
\newcommand{\PhaseThreeEnhaCotCometkiwi}{0.563}
\newcommand{\PhaseThreeEnhaCotComet}{0.698}
\newcommand{\PhaseThreeEnhaCotChrf}{35.9}
\newcommand{\PhaseThreeEnhaToolCometkiwi}{0.602}
\newcommand{\PhaseThreeEnhaToolComet}{0.78}
\newcommand{\PhaseThreeEnhaToolChrf}{45.5}
\newcommand{\PhaseThreeEnhaCottoolCometkiwi}{0.576}
\newcommand{\PhaseThreeEnhaCottoolComet}{0.736}
\newcommand{\PhaseThreeEnhaCottoolChrf}{40.5}
\newcommand{\PhaseThreeGapCottoolVsCotEndeCometkiwi}{-0.046}
\newcommand{\PhaseThreeGapCottoolVsCotEnhaCometkiwiLo}{0.013}
\newcommand{\PhaseThreeGapCottoolVsCotEnhaCometkiwiHi}{0.038}
\newcommand{\PhaseThreeGapCottoolVsToolEndeCometkiwi}{-0.099}
\newcommand{\PhaseThreeGapCottoolVsToolEnhaCometkiwi}{-0.027}
\newcommand{\PhaseThreeGapTestPLo}{0.004}
\newcommand{\PhaseThreeGapTestPHi}{0.024}
\newcommand{\PhaseThreeCotVerbatimRepeatEnde}{80}
\newcommand{\PhaseThreeCotVerbatimRepeatEnha}{72}
\newcommand{\PhaseThreeCottoolVerbatimRepeatEnde}{78}
\newcommand{\PhaseThreeCottoolVerbatimRepeatEnha}{86}
\newcommand{\PhaseThreeVerbatimRepeatDenominator}{100}
\newcommand{\PhaseThreeBestofnN}{5}
\newcommand{\PhaseThreeBestofnEnhaCometkiwi}{0.666}
\newcommand{\PhaseThreeBestofnEnhaSingleSampleCometkiwi}{0.581}

% --- phase4 ---
\newcommand{\PhaseFourModel}{Qwen2.5-32B-Instruct}
\newcommand{\PhaseFourNSentences}{100}
\newcommand{\PhaseFourNIterations}{4}
\newcommand{\PhaseFourCondAIterZeroCometkiwi}{0.81}
\newcommand{\PhaseFourCondAIterOneCometkiwi}{0.822}
\newcommand{\PhaseFourCondAIterTwoCometkiwi}{0.832}
\newcommand{\PhaseFourCondAIterThreeCometkiwi}{0.828}
\newcommand{\PhaseFourCondAIterFourCometkiwi}{0.832}
\newcommand{\PhaseFourCondBIterZeroCometkiwi}{0.81}
\newcommand{\PhaseFourCondBIterOneCometkiwi}{0.809}
\newcommand{\PhaseFourCondBIterTwoCometkiwi}{0.818}
\newcommand{\PhaseFourCondBIterThreeCometkiwi}{0.809}
\newcommand{\PhaseFourCondBIterFourCometkiwi}{0.816}
\newcommand{\PhaseFourCondAGainIterOneToPeakCometkiwi}{0.011}
\newcommand{\PhaseFourCondAPeakIteration}{4}
\newcommand{\PhaseFourCondBGainIterOneToPeakCometkiwi}{0.01}
\newcommand{\PhaseFourCondBPeakIteration}{2}
\newcommand{\PhaseFourCondAFrozenByIterFour}{31}
\newcommand{\PhaseFourCondBFrozenByIterThree}{48}
\newcommand{\PhaseFourNTotal}{100}
\newcommand{\PhaseFourRefineBugFirstRunIterZeroCometkiwi}{0.71}
\newcommand{\PhaseFourRefineBugFirstRunIterOneCometkiwi}{0.53}
\newcommand{\PhaseFourRefineBugFixedIterZeroCometkiwi}{0.81}
\newcommand{\PhaseFourExtractFixSuccessCount}{287}
\newcommand{\PhaseFourExtractFixTotalCount}{290}

% --- phase5 ---
\newcommand{\PhaseFiveBloomzModel}{bigscience/bloomz-7b1}
\newcommand{\PhaseFiveInkubalmModel}{lelapa/InkubaLM-0.4B}
\newcommand{\PhaseFiveInkubalmParamsB}{0.4}
\newcommand{\PhaseFiveNBugsFixedBloomzInkubalm}{9}
\newcommand{\PhaseFiveAyaTwoThreeExpanseNLanguages}{23}

% --- phase6 ---
\newcommand{\PhaseSixEnfrRootsBytes}{208242620434}
\newcommand{\PhaseSixEnswRootsBytes}{236482543}
\newcommand{\PhaseSixEnyoRootsBytes}{89695835}
\newcommand{\PhaseSixEnxhRootsBytes}{14304074}
\newcommand{\PhaseSixEnzuRootsBytes}{8511561}
\newcommand{\PhaseSixBloomDeRootsBytes}{0}
\newcommand{\PhaseSixEnfrCleanN}{10}
\newcommand{\PhaseSixEnfrTotalN}{10}
\newcommand{\PhaseSixEnswCleanN}{5}
\newcommand{\PhaseSixEnswTotalN}{10}
\newcommand{\PhaseSixEnyoCleanN}{0}
\newcommand{\PhaseSixEnyoTotalN}{10}
\newcommand{\PhaseSixEnxhCleanN}{0}
\newcommand{\PhaseSixEnxhTotalN}{10}
\newcommand{\PhaseSixEnzuCleanN}{0}
\newcommand{\PhaseSixEnzuTotalN}{10}
\newcommand{\PhaseSixEnfrCleanRatePct}{100}
\newcommand{\PhaseSixEnswCleanRatePct}{50}
\newcommand{\PhaseSixEnyoCleanRatePct}{0}
\newcommand{\PhaseSixEnxhCleanRatePct}{0}
\newcommand{\PhaseSixEnzuCleanRatePct}{0}
\newcommand{\PhaseSixMtFiveDeTokensB}{347}
\newcommand{\PhaseSixMtFiveHaTokensB}{0.2}
\newcommand{\PhaseSixMtFiveDePct}{3.05}
\newcommand{\PhaseSixMtFiveHaPct}{0.33}
\newcommand{\PhaseSixMtFiveDeHaRatio}{1735}



\setlist{noitemsep, topsep=2pt, parsep=0pt, partopsep=0pt}
\titlespacing*{\section}{0pt}{8pt}{4pt}
\titlespacing*{\subsection}{0pt}{6pt}{3pt}

\title{Agentic MT: Phases 1--6 Update}
\author{}
\date{}

\begin{document}
\maketitle
\vspace{-2em}

\section*{Where the theory stands}

Six phases in, the central joint-necessity claim (slow thinking and tool use are each
insufficient alone, jointly sufficient, and most valuable for low-resource pairs) has not
been confirmed or falsified --- every attempt to test it directly has run into a model
confound before producing a clean signal. What we do have is real, independent evidence
against the simplest instantiations of two supporting mechanisms (QE-in-the-loop, source-side
triggering), and a new, well-evidenced finding that a base model's raw pretraining exposure to
a language is itself a hard gate on generation, separate from anything the agentic scaffolding
does. Net effect: two of the five proposed research directions have real evidence against their
first-pass implementations, the flagship joint-necessity test remains open, and we've learned
something the theory didn't originally account for.

\section*{Findings, by implication}

\subsection*{Challenges to specific proposed mechanisms}

\begin{itemize}
  \item \textbf{QE-in-the-loop (Sec.~5.3) is not viable as specified for most of the target
    languages.} COMET-KIWI vs.\ human DA Spearman $r$ falls below the \PhaseOneFlagThreshold
    threshold for \PhaseOneNPairsFlagged of \PhaseOneNPairsTotal low-resource pairs tested:
    ha $=\PhaseOneHaSpearman$ (95\% CI [\PhaseOneHaCiLo, \PhaseOneHaCiHi], $n=\PhaseOneHaN$),
    ps $=\PhaseOnePsSpearman$, zu $=\PhaseOneZuSpearman$ --- versus an en-de benchmark of
    $\approx\PhaseOneEndeBenchmarkSpearman$. To trust this more: extend beyond the five pairs
    with real WMT DA/MQM human judgments; current coverage is a hard ceiling on what we can
    check.
  \item \textbf{Source-side triggering (Sec.~5.1) shows no signal from cheap features.} Five
    feature families (length, OOV proxy, LLM-rated idiomaticity, parse depth, domain markers)
    predict low quality at AUC $\PhaseTwoModelUniversalAuc$--$\PhaseTwoModelFullAuc$ against a
    chance baseline of $\PhaseTwoAucChance$. One promising-looking pooled correlation (OOV
    rate, pooled Pearson $r=\PhaseTwoOovPooledPearson$) evaporates to
    $\PhaseTwoOovPerlangLo$--$\PhaseTwoOovPerlangHi$ once controlled for language --- a genuine
    Simpson's paradox, not noise. To trust this more: test a real domain classifier and a
    frequency-list OOV measure (not a tokenizer-fragmentation proxy) before concluding
    triggering is a dead end rather than an implementation gap.
\end{itemize}

\subsection*{The central joint-necessity claim: confounded, not resolved}

\begin{itemize}
  \item \textbf{Phase 3} (Aya-101, en-de/en-ha, $n=\PhaseThreeNSentencesPerPair$/pair):
    tool-alone beat every other condition on both pairs, and the CoT-vs-tool gap differed
    significantly by resource level (paired-bootstrap $p=\PhaseThreeGapTestPLo$--$\PhaseThreeGapTestPHi$
    across metrics). But Aya-101 is not conversation-tuned, and its CoT chain produced
    byte-identical draft/refine/proofread output in
    \PhaseThreeCotVerbatimRepeatEnde--\PhaseThreeCottoolVerbatimRepeatEnha out of
    \PhaseThreeVerbatimRepeatDenominator sentences depending on pair/condition. The
    ``CoT hurts'' result is real in the sense that it's what happened, but is confounded with
    architecture, not a clean test of reasoning per se.
  \item \textbf{Phase 4} (Qwen2.5-32B, en-de only, $n=\PhaseFourNSentences$,
    \PhaseFourNIterations iterations): neither self-critique (condition A) nor external-QE
    refinement (condition B) peaked early and declined; both improved gradually to iteration
    \PhaseFourCondAPeakIteration/\PhaseFourCondBPeakIteration (gain
    $\PhaseFourCondAGainIterOneToPeakCometkiwi$/$\PhaseFourCondBGainIterOneToPeakCometkiwi$
    COMET-KIWI from iteration 1 to peak). This matches what Feng et al.\ actually report for
    TEaR, not the early-peak hypothesis we set out to test (that pattern belongs to a
    different paper). Secondary finding worth a follow-up: condition A froze only
    \PhaseFourCondAFrozenByIterFour/\PhaseFourNTotal sentences by iteration 4 vs.\ condition B's
    \PhaseFourCondBFrozenByIterThree/\PhaseFourNTotal by iteration 3, and A's trajectory is smoother
    --- span-level feedback (even self-generated) may matter more than feedback source
    (self vs.\ external). To trust this more: rerun condition B with XCOMET (span-level)
    instead of scalar COMET-KIWI to isolate granularity from source.
  \item \textbf{Phases 5--6}: two candidate replacement models for the low-resource arm were
    ruled out with direct evidence, not assumption. InkubaLM (\PhaseFiveInkubalmParamsB B
    params) is a capability floor --- incoherent even in English regardless of prompting.
    BLOOMZ needed \PhaseFiveNBugsFixedBloomzInkubalm mechanical fixes to reach a fair test,
    and then failed on resource-exposure grounds, not reliability grounds (see Phase 6 below).
    \textbf{Net: the joint-necessity question is still untested on a model without a
    disqualifying confound.}
\end{itemize}

\subsection*{Tangential but new: a pretraining-exposure floor}

\begin{itemize}
  \item \textbf{Phase 6}: BLOOMZ's clean-generation rate (l1\_baseline, manually categorized,
    $n=10$/pair) tracks its ROOTS pretraining byte count monotonically across four orders of
    magnitude: en-fr ($\PhaseSixEnfrRootsBytes$ bytes) $\to \PhaseSixEnfrCleanRatePct\%$ clean;
    en-sw ($\PhaseSixEnswRootsBytes$ bytes) $\to \PhaseSixEnswCleanRatePct\%$; en-xh
    ($\PhaseSixEnxhRootsBytes$ bytes) $\to \PhaseSixEnxhCleanRatePct\%$; en-zu
    ($\PhaseSixEnzuRootsBytes$ bytes) $\to \PhaseSixEnzuCleanRatePct\%$. Below some exposure
    level the model doesn't attempt the target language at all --- it defaults to echoing the
    English source (comprehension is intact; generation is not) --- rather than degrading
    gracefully. To trust this more: only 4 points on the curve; need intermediate-exposure
    languages (e.g.\ Yoruba, $\sim$90M bytes) to call it an actual regression rather than an
    eyeballed trend.
  \item Cross-checked the same question for Phase 3's Aya-101/mT5: German has
    $\PhaseSixMtFiveDeTokensB$B tokens in mC4 vs.\ Hausa's $\PhaseSixMtFiveHaTokensB$B
    ($\sim\PhaseSixMtFiveDeHaRatio\times$) --- so, unlike BLOOM's zero-exposure German, Phase 3's
    en-de anchor is legitimate for the model actually used there. Same question, opposite answer
    for the two models --- worth checking per-model going forward, not assuming either way.
\end{itemize}

\section*{Open decision}

Three ways to get an unconfounded joint-necessity result:

\begin{enumerate}
  \item \textbf{Keep searching for an open model that is both conversation-tuned and covers
    the target languages.} Two candidates already ruled out with evidence
    (\PhaseFiveNBugsFixedBloomzInkubalm fixes deep, in BLOOMZ's case). No known third candidate
    identified yet; another search cycle risks another multi-fix debugging odyssey with no
    guaranteed payoff, but preserves an all-open-source stack.
  \item \textbf{Split the two roles}: an API model for reasoning/orchestration, a
    Church-fine-tuned NLLB model queried as the tool, rather than asking one model to both
    translate and reason. This sidesteps Phase 6's exposure-floor problem entirely (NLLB is
    purpose-built and trained at scale on exactly these low-resource pairs, unlike a
    general-purpose LLM's incidental exposure), and arguably tests the theory's actual claim
    (reasoning \emph{and} tool use, not one model doing both) more faithfully than the current
    single-model design. Costs real API spend and departs from the fully-open constraint.
  \item \textbf{Narrow scope}: run the clean joint-necessity test on one pair BLOOMZ can
    already generate reliably (en-fr, $\PhaseSixEnfrCleanRatePct\%$ clean per Phase 6), and
    treat low-resource generalization as explicit future work rather than this paper's claim.
\end{enumerate}

\textbf{Leaning toward option 2} --- it resolves Phase 6's confound directly rather than
continuing to search for a single model that clears it. Option 3 is the fallback if preserving
an open-only stack matters more than resolving the question now.

\section*{What changes in the paper draft}

\begin{itemize}
  \item \textbf{Sec.~5.3 (QE-in-the-loop)}: needs an explicit caveat --- the proposed mechanism
    doesn't work for \PhaseOneNPairsFlagged/\PhaseOneNPairsTotal tested low-resource pairs;
    reframe as an open problem, not a solved detail.
  \item \textbf{Sec.~5.1 (triggering)}: same treatment --- cheap source-side features show no
    signal; needs either a caveat or a note that only crude implementations have been ruled
    out.
  \item \textbf{Sec.~1 (motivation)}: add the pretraining-exposure-floor finding (Phase 6) as a
    distinct failure mode from reasoning/tool insufficiency --- ``low-resource'' can mean
    ``the base model never learned to generate this language at all,'' independent of any
    agentic scaffolding.
  \item \textbf{Core joint-necessity claim}, wherever stated: needs a
    ``confounded/untested cleanly'' framing --- not silence, not overclaiming from Phase 3's
    suggestive-but-confounded numbers.
\end{itemize}

\appendix
\section*{Appendix: debugging record (compressed)}

\begin{itemize}
  \item \textbf{Aya-101 (Phase 3)}: T5 seq2seq, no chat template; multi-turn context handled by
    flattening to a role-tagged transcript. Worked for single-turn conditions; degenerated for
    the 4-step CoT chain (see verbatim-repeat numbers above).
  \item \textbf{Refine-step scoring bug (Phase 4)}: Qwen2.5-32B's refine responses were often
    full English explanations (``Given the feedback\ldots the revised sentence would
    be:\ldots''), and the entire response was being scored as the translation --- produced a
    fake collapse from $\PhaseFourRefineBugFirstRunIterZeroCometkiwi \to \PhaseFourRefineBugFirstRunIterOneCometkiwi$
    COMET-KIWI. Fixed with an explicit ``output only the translation'' instruction plus a
    post-hoc extractor, validated against the original buggy data
    ($\PhaseFourExtractFixSuccessCount$/$\PhaseFourExtractFixTotalCount$ cleaned correctly)
    before re-running.
  \item \textbf{BLOOMZ/InkubaLM (Phase 5)}: \PhaseFiveNBugsFixedBloomzInkubalm fixes in
    sequence --- missing pre-download, an outdated KV-cache code path in InkubaLM's custom
    model class, partial gated-repo caching, a context-window miscalculation, role-tag framing
    that made both models echo prompt boilerplate instead of translating, a completion-cue
    removal that caused empty output, a misdiagnosed repetition-penalty fix, a cue-induced
    short-answer truncation bug, and finally a genuine reliability problem (inconsistent
    correct-translation-vs.-source-echo behavior) that was not further fixable. An earlier full
    run that looked like an interesting contradiction of Phase 3 was explicitly retracted once
    these bugs were understood --- flagged as uncitable in \texttt{phase\_5\_summary.md} rather
    than left standing.
  \item \textbf{ROOTS-table correction (Phase 6)}: a secondary characterization claimed
    Xhosa/Zulu were entirely absent from BLOOM's pretraining data. Checked the primary source
    (BLOOM paper, Table 1) directly --- both are present with real, if small, byte counts
    ($\PhaseSixEnxhRootsBytes$ / $\PhaseSixEnzuRootsBytes$ bytes respectively). The claim that
    German has zero presence, however, held up on the same check --- and turned out to be the
    more consequential finding (Phase 3's/Phase 5's en-de anchor question, resolved oppositely
    for the two models involved).
\end{itemize}

\end{document}
